\begin{annexitem}{annex:p3.hyp.1}{thm:p03-hyp-projective}{Presenting truth projectively}
\begin{proof}
Work first in a finite relational signature \(\rho\). For a
\(\rho\)-structure \(M\), use its objects as urelements of an
auxiliary model of KPU. The finite presentation in
\cref{thm:p03-recursive-auxiliary-presentation} specifies a named
nonempty domain \(S\), a membership relation \(E\), a set of all
its urelements, and standard internal \(\omega\). It is available
in either input logic by \cref{thm:p03-standard-numbers}.
Identify the distinguished reduct domain with exactly those
urelements. This identification is a first-order condition in the
ambient expanded signature; the reduct domain may be empty.

Add finitely many named elements of \(S\) coding the relations of
\(M\). For an \(n\)-ary relation require that the members of its
code are exactly the internal \(n\)-tuple codes of urelements whose
external tuple satisfies that relation. The fixed finite tuple
construction is first-order definable from membership, and KPU proves
its existence and uniqueness. Require also that every member of the
relation code is such a tuple, so no hidden members remain. At arity
zero the unique empty tuple is included precisely when the nullary
relation is true. Code the resulting proper structure by the domain
set and this finite list of relation codes; call its code \(m\).
All these are finitely many first-order requirements.

Fix the particular formula code \(\varphi\in\mathrm{HYP}\).
\Cref{thm:p03-hyp-fixed-names} gives its parameter-free
\(\Sigma_1\) name \(\sigma_\varphi(x)\).
\Cref{thm:p03-adm-satisfaction} supplies the two complementary
\(\Sigma_1\) evaluation formulas
\(\operatorname{Sat}(m,x)\) and
\(\operatorname{NotSat}(m,x)\) for sentences, using the empty
assignment. Conjoin respectively
\[
 \exists x\in S\,
   \bigl(\sigma_\varphi^S(x)\land\operatorname{Sat}^S(m,x)\bigr),
 \qquad
 \exists x\in S\,
   \bigl(\sigma_\varphi^S(x)\land\operatorname{NotSat}^S(m,x)\bigr).
\]
Here the superscript means ordinary first-order relativization to the
set-theoretic structure \((S,E)\), with its urelement predicate.
These are finite formulas, even when \(\varphi\) is infinitary:
\(\sigma_\varphi\) names its code rather than spelling out its
constituents in the input logic. With the preceding auxiliary-model
sentence they give the two proposed projective presentations.

To verify soundness, take any witnessing auxiliary model \(B\),
possibly externally ill founded. By
\cref{thm:p03-adm-truncation}, its well-founded part collapses to a
transitive admissible set over the same urelements, with the actual
\(\omega\) as an element. Denote this part by \(B_0\), identifying
it with its collapse. It contains \(\mathrm{HYP}\) by
\cref{thm:p03-hyp-least}. It also contains the code \(m\): the set
of all urelements is well founded, every finite tuple of urelements
is well founded, and each relation code has only such tuples as
members. The finite structure record is therefore well founded.
Its collapse codes exactly the distinguished reduct \(M\).

The fixed-name theorem says that the object satisfying
\(\sigma_\varphi\) in \(B\) is the actual old code \(\varphi\).
In \(B_0\), the evaluation of that code on \(m\) exists and is
its actual truth value in \(M\). Its positive or negative
\(\Sigma_1\) certificate persists to the end extension \(B\).
KPU proves the two evaluations complementary on legal inputs, so
\(B\) cannot assert the opposite value. Consequently the positive
presentation implies \(M\models\varphi\), and the negative one
implies \(M\not\models\varphi\).

For existence, choose in the surrounding set theory a regular
uncountable \(\kappa\) larger than the size of the domain and all
its finitely many relations. The standard admissible construction
\(H(\kappa)\) over these urelements contains their entire domain
set, their relation codes, the finite structure record, and
\(\omega\). The proof of
\cref{thm:p03-w03-p03-039-nadel-hkappa} applies with the given
urelements: bounded Separation preserves hereditary size below
\(\kappa\), and regularity collects fewer than \(\kappa\)
witnesses of hereditary size below \(\kappa\). It is therefore
admissible and contains \(\mathrm{HYP}\). Its correct evaluation
of \(\varphi\) has the required positive or negative certificate.
Expand it by the external arithmetic, assignment codes, and truth
relation from \cref{thm:p03-recursive-auxiliary-presentation}; its
internal natural numbers are standard. This gives a witness for the
appropriate presentation. No countability assumption on \(M\) was
used. The same argument works for an empty urelement domain.

For functions and constants, use total functional relation graphs,
whose codes are handled in exactly the same way. The standard graph
translation preserves truth of \(\varphi\), and first-order graph
conditions recover the original signature. This proves the claim for
every finite signature.
\end{proof}
\end{annexitem}

\begin{annexitem}{annex:p3.hyp.2}{thm:p03-hyp-characterization}{Equality of expressive power}
\begin{proof}
First both input logics translate into \(L_{\mathrm{HYP}}\).
Replace \(Q_0x\,\psi(x)\) by the conjunction, over \(n<\omega\),
asserting the existence of \(n\) distinct objects each satisfying
\(\psi\). In weak second-order logic replace an existential finite
set quantifier by the disjunction over \(n<\omega\) of the formulas
obtained by enumerating that set with \(n\) first-order variables.
Replace an occurrence of membership in the quantified set by the
finite disjunction of equalities to those variables. Repetitions in
the enumeration are harmless, and the \(n=0\) disjunct treats the
empty set. Universal finite-set quantification is obtained by
negation. If finite relations of arity greater than one are included
in the weak second-order syntax, enumerate finitely many tuples
instead, and replace a relation atom by the disjunction of the
coordinatewise equalities to those tuples. Fresh-variable choices make
substitution capture free.

Induction on the finite input formula proves these translations
correct. Each countable family of clauses is the range of a
\(\Sigma\)-definable function on the actual set \(\omega\) in
\(\mathrm{HYP}\), with the previously translated constituents as
parameters. \Cref{thm:p03-adm-collection,thm:p03-adm-syntax} therefore
places the family and its conjunction or disjunction in
\(\mathrm{HYP}\). The translated formula has only the original
finite signature and its original free variables. This proves
\(L_w,L(Q_0)\leq L_{\mathrm{HYP}}\), including empty structures.

By \cref{thm:p03-adm-relativized-separation},
\(L_{\mathrm{HYP}}\) has the Souslin--Kleene property for the
relativized projective convention. Apply the leastness clause of
\cref{thm:p03-delta-closure} to the preceding two extensions. It gives
\[
 \Delta(L_w)\leq L_{\mathrm{HYP}},\qquad
 \Delta(L(Q_0))\leq L_{\mathrm{HYP}}.
\]
Conversely \cref{thm:p03-hyp-projective} gives complementary projective
presentations, in each input logic, for every
\(L_{\mathrm{HYP}}\)-sentence. Part~3 of
\cref{thm:p03-delta-closure} therefore makes each such sentence
definable in each closure. This proves both reverse inclusions and
the displayed equality. The leastness assertion is precisely the
closure theorem's leastness assertion, now with its extension
identified explicitly.
\end{proof}
\end{annexitem}
